% ============================================
% DOC-03 — Laudo de Necropsia (edição revista R1075)
% ============================================
\fragtitle{doc}{D}{DOC}{03}{Laudo de Necropsia}
\markboth{DOC-03 --- Laudo de Necropsia}{}

\noindent \textit{Instituto Médico Legal de Barbacena --- 14 de junho de 1979 --- Dr. Álvaro Torres}

\noindent \textbf{Paciente:} Joaquim Antônio Correia \quad \textbf{Registro:} 1.260 \quad \textbf{Idade estimada:} 72 anos \quad \textbf{Óbito:} 12 de junho de 1979 \quad \textbf{Necropsia:} 14 de junho de 1979 \quad \textbf{Legista:} Dr. Álvaro Torres --- CRM-MG 3.117.

\subsection*{1. Exame externo}

\noindent \textbf{1.1 Corpo.} Peso 42 kg; estatura 1,62 m; índice de massa corporal 16,0. Caquexia acentuada. Panículo adiposo ausente; saliências ósseas visíveis sob a pele. Pele ressecada, com pregueamento persistente. Vinte e um parasitas \textit{Ascaris lumbricoides} foram contados no intestino; a carga parasitária competia com o corpo pelos poucos nutrientes disponíveis.

\noindent \textbf{1.2 Olho direito.} A esclera direita apresenta coloração amarelo-ouro intensa, homogênea, da córnea à conjuntiva bulbar. A esclera esquerda é branca, com vascularização normal. A coloração não é compatível com icterícia hepática clássica, que seria bilateral. À palpação, o globo direito é ligeiramente mais firme. Não há trauma recente, prótese ou corpo estranho. A mancha não desbota à compressão.

\noindent \textit{Anotação marginal, caligrafia diferente:} ``Depois da exumação, esta hipótese caiu por terra. O olho brilhava no escuro. Não há explicação médica para o fenômeno. --- Lúcia M.''

\noindent \textbf{1.3 Eletroconvulsoterapia.} Marcas ovais de queimadura elétrica nas regiões temporais, bilaterais, 3 cm × 2 cm, com alopecia permanente. Fratura por compressão de T5, consolidada; fratura antiga do colo do fêmur esquerdo, com encurtamento de 2 cm; fratura consolidada do ramo mandibular direito. As lesões são compatíveis com 37 sessões de eletrochoque sem anestesia, sem gel condutor e sem proteção bucal, entre 1937 e 1952.

\noindent \textbf{1.4 Lobotomia transorbital.} Cicatrizes puntiformes bilaterais nos cantos internos dos olhos, 2 mm. Na lâmina cribriforme, dois orifícios de 4 mm, simétricos, a 2,5 cm da linha média, com bordas lisas e cicatrizadas. Os orifícios estão pérvios. O prontuário não contém autorização para o procedimento.

\noindent \textbf{1.5 Contenção e automutilação.} Faixas de hiperpigmentação circumferencial nos punhos e tornozelos, compatíveis com contenção prolongada. Contratura dos dedos em semiflexão. Escoriações lineares no tórax e no abdome. Cicatrizes de arranhões profundos nos antebraços, com terra impregnada na derme. Amputação parcial da falange distal do dedo mínimo esquerdo, sem registro de trauma.

\noindent \textbf{1.6 Queimaduras.} Manchas de hipopigmentação no tórax e no abdome, 2 a 8 cm, compatíveis com escaldadura por imersão nos chamados ``banhos terapêuticos''.

\subsection*{2. Exame interno}

\noindent \textbf{2.1 Causa mortis.} Parada cardiorrespiratória secundária a insuficiência cardíaca congestiva, caquexia e broncoaspiração. Coração de 220 g, com atrofia marrom do miocárdio. Estômago vazio; mucosa atrófica. Fígado de 980 g, com esteatose moderada. Não há cirrose.

\noindent \textbf{2.2 Encéfalo.} Peso 1.120 g. Hemisférios simétricos; meninges transparentes. Não há gliose, cavitação ou alteração cortical compatível com lobotomia. As áreas orbitofrontais e pré-frontais foram dissecadas. A via de acesso dos leucótomos existe; o parênquima, não foi atingido de forma detectável.

\noindent \textit{Nota do legista:} ``A ausência de dano cerebral não é um alívio. É o achado mais perturbador. Joaquim não era violento por lesão cerebral. A violência era dele. Ou era de outra coisa. Os procedimentos não deixaram marca no cérebro; deixaram no corpo inteiro. O cérebro calado é a parte mais gritante do laudo.''

\noindent \textbf{2.3 Coluna.} Fraturas por compressão em T3, T5 e T8, estáveis, sem déficit neurológico. Tríade compatível com convulsões tônico-clônicas repetidas.

\subsection*{3. Síntese}

\begin{itemize}
\item ECT sem anestesia: queimaduras temporais, fraturas vertebrais, fratura mandibular e femoral, perda dentária.
\item Lobotomia transorbital: perfuração bilateral da lâmina cribriforme, sem lesão cortical detectável.
\item Contenção prolongada: faixas em punhos e tornozelos, contratura digital, fibrose em coxas.
\item Banhos terapêuticos: cicatrizes térmicas no tórax e no abdome.
\end{itemize}

\noindent \textbf{Causa mortis:} parada cardiorrespiratória. \textbf{Causa contribuinte:} caquexia por desnutrição crônica. \textbf{Comorbidade:} alta carga parasitária. Não foram encontrados marcadores orgânicos que justifiquem 62 anos de institucionalização. A lesão mais significativa do paciente não é anatômica.

\vspace{0.5em}
\noindent Dr. Álvaro Torres --- Médico Legista --- IML Barbacena --- CRM-MG 3.117

\vspace{0.8em}
\noindent \textit{Adendo manuscrito, aparentemente posterior:} ``Este laudo foi redigido antes da exumação. Os achados da seção seguinte foram acrescentados após a abertura da cova e a recuperação do diário. O signatário declara que não há explicação médica para os fenômenos descritos no item 4, limitando-se a relatar o que foi apresentado pelo coveiro e pelo material recuperado. O diário permanece sob custódia do IML, anexado a este processo. --- A.T.''

\vspace{0.5em}
\noindent \textit{Nota adicional, caligrafia de L.M., protocolo 2026:} ``Fragmentos ósseos de Strigiformes (coruja-buraqueira) foram recuperados do mesmo estrato da cova. A presença de restos de coruja no sepultamento de Joaquim foi registrada pelo coveiro como \textit{a ave que não saiu de perto}. Não foi encontrada explicação tafonômica para a deposição. Ossos de coruja não são achado comum em sepultamentos humanos no Colônia.''

\noindent\textit{\textbf{↪ Rotas:} \rota{MEM-04} \rotasep \rota{DOC-02} \rotasep \rota{LUC-03} \rotasep \rota{CONT-03}}
